% GENERATED FILE. Edit canonical lessons, then run npm run generate:books.
% canonical-source-hash: fnv1a64:3e6ba2a8e87f73c6
% canonical-lessons: SA-C404-prantaram, SA-C404-kantarah, SA-C404-desatanam, SA-C404-niliyate, SA-C404-sunyah

\chapter{Lonely Roads and Wild Places}
\label{ch:sa-lonely-roads-and-wild-places}
\hlchaptermodality{404}

\begin{chapteropening}
\textbf{By the end of this chapter:} \emph{I can talk about a wild, empty stretch and where to take cover.}

Complete the last lesson of chapter 404: I can talk about a wild, empty stretch and where to take cover.
\end{chapteropening}

\section[\texorpdfstring{\sk{प्रान्तरम्} (prāntaram)}{प्रान्तरम् (prāntaram)}]{\sk{प्रान्तरम्} (prāntaram) — a long lonely road}

\label{lesson:SA-C404-prantaram}



\begin{quote}
Before the new one: say the Sanskrit for a traveller on foot, then the Sanskrit for roams.
\end{quote}

\subsection*{You'll want to know: \sk{प्रान्तरम्}}
\textbf{\sk{प्रान्तरम्}} — \emph{prāntaram} — \textquotedblleft{}a long lonely road\textquotedblright{}.

\begin{quote}
\sk{प्रान्तरे} \sk{कः} \sk{अपि} \sk{न} \sk{अस्ति।}
\end{quote}

\begin{quote}
\emph{prāntare kaḥ api na asti.}
\end{quote}

\begin{quote}
There is nobody at all on the lonely road.
\end{quote}

A neuter noun.

\begin{cousinweb}
\textbf{\sk{प्रान्तरम्}} (\emph{prāntaram}) is built from \emph{pra}, \textquotedblleft{}forth\textquotedblright{}, and \emph{antara}, \textquotedblleft{}an interval, the space between\textquotedblright{}; MW gives \emph{prāntara} as a long desolate road and as the country between two villages (MW).
\end{cousinweb}

\begin{grammarlens}[title={What you've built}]
One more word for getting around.
\end{grammarlens}

\subsection*{Your turn}
\begin{itemize}
\raggedright
  \item \emph{Say it:} \emph{prāntaram}
  \item \emph{Say it:} \emph{prāntaram}, once more
  \item \emph{Say it:} \emph{prāntare kaḥ api na asti.} — There is nobody at all on the lonely road.
  \item \emph{Recall:} say the Sanskrit for a traveller on foot, then the Sanskrit for roams, then say \emph{prāntaram} again
\end{itemize}

\begin{tcolorbox}[breakable,colback=teal!4,colframe=teal!35!black,title={Before you move on}]
What does \sk{प्रान्तरम्} mean? (\textquotedblleft{}A long lonely road\textquotedblright{}.) Say it once more.
\end{tcolorbox}
\section[\texorpdfstring{\sk{कान्तारः} (kāntāraḥ)}{कान्तारः (kāntāraḥ)}]{\sk{कान्तारः} (kāntāraḥ) — a wilderness; a forest track}

\label{lesson:SA-C404-kantarah}



\begin{quote}
Before the new one: say the Sanskrit for roams, then the Sanskrit for a long lonely road.
\end{quote}

\subsection*{You'll want to know: \sk{कान्तारः}}
\textbf{\sk{कान्तारः}} — \emph{kāntāraḥ} — \textquotedblleft{}a wilderness; a forest track\textquotedblright{}.

\begin{quote}
\sk{कान्तारे} \sk{सिंहाः} \sk{वृकाः} \sk{च} \sk{वसन्ति।}
\end{quote}

\begin{quote}
\emph{kāntāre siṃhāḥ vṛkāḥ ca vasanti.}
\end{quote}

\begin{quote}
Lions and wolves live in the wilderness.
\end{quote}

A masculine noun.

\begin{cousinweb}
\textbf{\sk{कान्तारः}} (\emph{kāntāraḥ}) is the word MW gives for a large wood or wilderness and for a difficult road through a forest; MW gives no derivation, so its origin is uncertain (MW).
\end{cousinweb}

\begin{grammarlens}[title={What you've built}]
One more word for getting around.
\end{grammarlens}

\subsection*{Your turn}
\begin{itemize}
\raggedright
  \item \emph{Say it:} \emph{kāntāraḥ}
  \item \emph{Say it:} \emph{kāntāraḥ}, once more
  \item \emph{Say it:} \emph{kāntāre siṃhāḥ vṛkāḥ ca vasanti.} — Lions and wolves live in the wilderness.
  \item \emph{Recall:} say the Sanskrit for roams, then the Sanskrit for a long lonely road, then say \emph{kāntāraḥ} again
\end{itemize}

\begin{tcolorbox}[breakable,colback=teal!4,colframe=teal!35!black,title={Before you move on}]
What does \sk{कान्तारः} mean? (\textquotedblleft{}A wilderness; a forest track\textquotedblright{}.) Say it once more.
\end{tcolorbox}
\section[\texorpdfstring{\sk{देशाटनम्} (deśāṭanam)}{देशाटनम् (deśāṭanam)}]{\sk{देशाटनम्} (deśāṭanam) — travel through the country}

\label{lesson:SA-C404-desatanam}



\begin{quote}
Before the new one: say the Sanskrit for a long lonely road, then the Sanskrit for a wilderness.
\end{quote}

\subsection*{You'll want to know: \sk{देशाटनम्}}
\textbf{\sk{देशाटनम्}} — \emph{deśāṭanam} — \textquotedblleft{}travel through the country\textquotedblright{}.

\begin{quote}
\sk{देशाटनेन} \sk{ज्ञानं} \sk{वर्धते।}
\end{quote}

\begin{quote}
\emph{deśāṭanena jñānaṃ vardhate.}
\end{quote}

\begin{quote}
Knowledge grows by travelling through the country.
\end{quote}

A neuter noun.

\begin{cousinweb}
\textbf{\sk{देशाटनम्}} (\emph{deśāṭanam}) is built from \emph{deśa}, \textquotedblleft{}a country\textquotedblright{}, and \emph{aṭana}, \textquotedblleft{}roaming\textquotedblright{}, from the root √aṭ, \textquotedblleft{}to roam\textquotedblright{}, which also gives \emph{paryaṭaka}, \textquotedblleft{}a tourist\textquotedblright{}; MW gives the compound as roaming through a land, travelling (MW).
\end{cousinweb}

\begin{grammarlens}[title={What you've built}]
One more word for getting around.
\end{grammarlens}

\subsection*{Your turn}
\begin{itemize}
\raggedright
  \item \emph{Say it:} \emph{deśāṭanam}
  \item \emph{Say it:} \emph{deśāṭanam}, once more
  \item \emph{Say it:} \emph{deśāṭanena jñānaṃ vardhate.} — Knowledge grows by travelling through the country.
  \item \emph{Recall:} say the Sanskrit for a long lonely road, then the Sanskrit for a wilderness, then say \emph{deśāṭanam} again
\end{itemize}

\begin{tcolorbox}[breakable,colback=teal!4,colframe=teal!35!black,title={Before you move on}]
What does \sk{देशाटनम्} mean? (\textquotedblleft{}Travel through the country\textquotedblright{}.) Say it once more.
\end{tcolorbox}
\section[\texorpdfstring{\sk{निलीयते} (nilīyate)}{निलीयते (nilīyate)}]{\sk{निलीयते} (nilīyate) — takes cover; hides oneself}

\label{lesson:SA-C404-niliyate}



\begin{quote}
Before the new one: say the Sanskrit for a wilderness, then the Sanskrit for travel through the country.
\end{quote}

\subsection*{You'll want to know: \sk{निलीयते}}
\textbf{\sk{निलीयते}} — \emph{nilīyate} — \textquotedblleft{}takes cover; hides oneself\textquotedblright{}.

\begin{quote}
\sk{पथिकः} \sk{सिंहात्} \sk{गुहायां} \sk{निलीयते।}
\end{quote}

\begin{quote}
\emph{pathikaḥ siṃhāt guhāyāṃ nilīyate.}
\end{quote}

\begin{quote}
The wayfarer takes cover from the lion in a cave.
\end{quote}

Past (he …): \sk{सः} \sk{गुहायां} \sk{निलीनवान्}, \emph{saḥ guhāyāṃ nilīnavān}.

\begin{cousinweb}
\textbf{\sk{निलीयते}} (\emph{nilīyate}) is built from \emph{ni}, \textquotedblleft{}down\textquotedblright{}, and the root √lī, \textquotedblleft{}to cling, to settle\textquotedblright{}; MW gives \emph{ni-lī} as to settle down or alight and to hide oneself from something (in the ablative), and KRM gives the participle \emph{līna} (MW).
\end{cousinweb}

\begin{grammarlens}[title={What you've built}]
One more word for getting around.
\end{grammarlens}

\subsection*{Your turn}
\begin{itemize}
\raggedright
  \item \emph{Say it:} \emph{nilīyate}
  \item \emph{Say it:} \emph{nilīyate}, once more
  \item \emph{Say it:} \emph{pathikaḥ siṃhāt guhāyāṃ nilīyate.} — The wayfarer takes cover from the lion in a cave.
  \item \emph{Recall:} say the Sanskrit for a wilderness, then the Sanskrit for travel through the country, then say \emph{nilīyate} again
\end{itemize}

\begin{tcolorbox}[breakable,colback=teal!4,colframe=teal!35!black,title={Before you move on}]
What does \sk{निलीयते} mean? (\textquotedblleft{}Takes cover; hides oneself\textquotedblright{}.) Say it once more.
\end{tcolorbox}
\section[\texorpdfstring{\sk{शून्यः} (śūnyaḥ)}{शून्यः (śūnyaḥ)}]{\sk{शून्यः} (śūnyaḥ) — deserted; empty}

\label{lesson:SA-C404-sunyah}



\begin{quote}
Before the new one: say the Sanskrit for travel through the country, then the Sanskrit for takes cover.
\end{quote}

\subsection*{You'll want to know: \sk{शून्यः}}
\textbf{\sk{शून्यः}} — \emph{śūnyaḥ} — \textquotedblleft{}deserted; empty\textquotedblright{}.

\begin{quote}
\sk{रात्रौ} \sk{मार्गः} \sk{शून्यः} \sk{भवति।}
\end{quote}

\begin{quote}
\emph{rātrau mārgaḥ śūnyaḥ bhavati.}
\end{quote}

\begin{quote}
At night the road is deserted.
\end{quote}

\begin{cousinweb}
\textbf{\sk{शून्यः}} (\emph{śūnyaḥ}) is the word MW gives for empty, void, deserted; MW gives no sure derivation, so its origin is uncertain (MW).
\end{cousinweb}

\begin{grammarlens}[title={What you've built}]
One more word for getting around.
\end{grammarlens}

\subsection*{Your turn}
\begin{itemize}
\raggedright
  \item \emph{Say it:} \emph{śūnyaḥ}
  \item \emph{Say it:} \emph{śūnyaḥ}, once more
  \item \emph{Say it:} \emph{rātrau mārgaḥ śūnyaḥ bhavati.} — At night the road is deserted.
  \item \emph{Recall:} say the Sanskrit for travel through the country, then the Sanskrit for takes cover, then say \emph{śūnyaḥ} again
\end{itemize}

\begin{tcolorbox}[breakable,colback=teal!4,colframe=teal!35!black,title={Before you move on}]
What does \sk{शून्यः} mean? (\textquotedblleft{}Deserted; empty\textquotedblright{}.) Say it once more.
\end{tcolorbox}
